% !TITLE 楚辞·离骚（节选）
% !AUTHOR 屈原
% !DYNASTY 先秦
% !GENRE 楚辞
% !ORDER 12
% !SOURCE 楚辞·离骚

\begin{poem}
帝高阳之苗裔兮\textsuperscript{1}，朕皇考曰伯庸\textsuperscript{2}。
摄提贞于孟陬兮\textsuperscript{3}，惟庚寅吾以降\textsuperscript{4}。
皇览揆余初度兮\textsuperscript{5}，肇锡余以嘉名\textsuperscript{6}：
名余曰正则兮\textsuperscript{7}，字余曰灵均\textsuperscript{8}。
纷吾既有此内美兮\textsuperscript{9}，又重之以修能\textsuperscript{10}。
扈江离与辟芷兮\textsuperscript{11}，纫秋兰以为佩\textsuperscript{12}。
汩余若将不及兮\textsuperscript{13}，恐年岁之不吾与\textsuperscript{14}。
朝搴阰之木兰兮\textsuperscript{15}，夕揽洲之宿莽\textsuperscript{16}。
日月忽其不淹兮\textsuperscript{17}，春与秋其代序\textsuperscript{18}。
惟草木之零落兮\textsuperscript{19}，恐美人之迟暮\textsuperscript{20}。
不抚壮而弃秽兮\textsuperscript{21}，何不改乎此度\textsuperscript{22}？
乘骐骥以驰骋兮\textsuperscript{23}，来吾道夫先路\textsuperscript{24}！

……

长太息以掩涕兮\textsuperscript{25}，哀民生之多艰\textsuperscript{26}。
余虽好修姱以鞿羁兮\textsuperscript{27}，謇朝谇而夕替\textsuperscript{28}。
既替余以蕙纕兮\textsuperscript{29}，又申之以揽茝\textsuperscript{30}。
亦余心之所善兮\textsuperscript{31}，虽九死其犹未悔\textsuperscript{32}。
怨灵修之浩荡兮\textsuperscript{33}，终不察夫民心\textsuperscript{34}。
众女嫉余之蛾眉兮\textsuperscript{35}，谣诼谓余以善淫\textsuperscript{36}。
固时俗之工巧兮\textsuperscript{37}，偭规矩而改错\textsuperscript{38}。
背绳墨以追曲兮\textsuperscript{39}，竞周容以为度\textsuperscript{40}。
忳郁邑余侘傺兮\textsuperscript{41}，吾独穷困乎此时也\textsuperscript{42}。
宁溘死以流亡兮\textsuperscript{43}，余不忍为此态也\textsuperscript{44}。
鸷鸟之不群兮\textsuperscript{45}，自前世而固然\textsuperscript{46}。
何方圜之能周兮\textsuperscript{47}，夫孰异道而相安\textsuperscript{48}？
屈心而抑志兮\textsuperscript{49}，忍尤而攘诟\textsuperscript{50}。
伏清白以死直兮\textsuperscript{51}，固前圣之所厚\textsuperscript{52}。

……

悔相道之不察兮\textsuperscript{53}，延伫乎吾将反\textsuperscript{54}。
回朕车以复路兮\textsuperscript{55}，及行迷之未远\textsuperscript{56}。
步余马于兰皋兮\textsuperscript{57}，驰椒丘且焉止息\textsuperscript{58}。
进不入以离尤兮\textsuperscript{59}，退将复修吾初服\textsuperscript{60}。
制芰荷以为衣兮\textsuperscript{61}，集芙蓉以为裳\textsuperscript{62}。
不吾知其亦已兮\textsuperscript{63}，苟余情其信芳\textsuperscript{64}。
高余冠之岌岌兮\textsuperscript{65}，长余佩之陆离\textsuperscript{66}。
芳与泽其杂糅兮\textsuperscript{67}，唯昭质其犹未亏\textsuperscript{68}。
忽反顾以游目兮\textsuperscript{69}，将往观乎四荒\textsuperscript{70}。
佩缤纷其繁饰兮\textsuperscript{71}，芳菲菲其弥章\textsuperscript{72}。
民生各有所乐兮\textsuperscript{73}，余独好修以为常\textsuperscript{74}。
虽体解吾犹未变兮\textsuperscript{75}，岂余心之可惩\textsuperscript{76}。
\end{poem}

\begin{annotations}
\notetext{1}{帝高阳之苗裔：我是古帝高阳（颛顼）的后代。苗裔（yì），后代子孙。屈原开篇先追溯祖先，表明自己出身高贵。}
\notetext{2}{朕皇考曰伯庸：我去世的父亲名叫伯庸。朕（zhèn），我（秦以前人人可自称朕）；皇考，对去世父亲的尊称。}
\notetext{3}{摄提贞于孟陬：摄提格（寅年）那年，正当孟春正月。摄提，岁星纪年法中的寅年；贞，正当；孟陬（zōu），孟春正月。}
\notetext{4}{惟庚寅吾以降：在庚寅日那天我降生了。庚寅，干支纪日。}
\notetext{5}{皇览揆余初度：父亲观察揣度我初生时的气度。览，观察；揆（kuí），揣度；初度，初生的气度。}
\notetext{6}{肇锡余以嘉名：开始赐给我美好的名字。肇（zhào），开始；锡，同"赐"。}
\notetext{7}{正则：公正而有法则。"正则"是屈原名"平"的寓意——平正可法。}
\notetext{8}{灵均：灵善而均平。"灵均"是屈字"原"的寓意——地之善而均平者。}
\notetext{9}{纷吾既有此内美：我既有这么多内在的美质。纷，盛多的样子。}
\notetext{10}{又重之以修能：又加上外在的杰出才能。重（chóng），加上；修能，美好的才能。}
\notetext{11}{扈江离与辟芷：披着江离和辟芷（香草）。扈（hù），披；江离，川芎；辟，同"僻"，幽僻；芷（zhǐ），白芷。}
\notetext{12}{纫秋兰以为佩：把秋兰编成绳索当作佩饰。纫（rèn），连缀、编织。}
\notetext{13}{汩余若将不及兮：时光如水流逝，我像是追赶不上。汩（yù），水流迅疾的样子。}
\notetext{14}{恐年岁之不吾与：担心岁月不等待我。不吾与，"不与吾"的倒装——不等我。}
\notetext{15}{朝搴阰之木兰：清晨去山上拔木兰。搴（qiān），拔取；阰（pí），山坡。}
\notetext{16}{夕揽洲之宿莽：黄昏去沙洲上采宿莽。宿莽，一种卷施草，冬天不死。朝夕采摘香草比喻努力修身。}
\notetext{17}{日月忽其不淹兮：日月匆匆，不肯停留。忽，快速；淹，停留。}
\notetext{18}{春与秋其代序：春天与秋天交替更迭。}
\notetext{19}{惟草木之零落兮：想到草木会凋零。惟，思、念。}
\notetext{20}{恐美人之迟暮：担心美人（比喻君主或自身理想）到了衰老之年。"美人"在楚辞中常喻楚王。}
\notetext{21}{不抚壮而弃秽兮：不趁着壮年去除污秽。抚，趁着；壮，壮年。}
\notetext{22}{何不改乎此度：为什么不改变这种态度（指腐败的政治风气）呢？}
\notetext{23}{乘骐骥以驰骋兮：骑着骏马奔驰。骐骥（qí jì），千里马。}
\notetext{24}{来吾道夫先路：来吧，让我在前面引路。道，同"导"；先路，前面的路。}
\notetext{25}{长太息以掩涕兮：长长地叹息，擦着眼泪。太息，叹息；掩涕，抹泪。}
\notetext{26}{哀民生之多艰：哀叹人生如此多难。民生，人生。}
\notetext{27}{余虽好修姱以鞿羁兮：我虽然喜好美德却受到束缚。修姱（kuā），美好；鞿羁（jī jī），马缰绳和马笼头，比喻受到牵制约束。}
\notetext{28}{謇朝谇而夕替：早上进谏，晚上就被罢免。謇（jiǎn），直言敢谏；谇（suì），谏言；替，废黜、罢免。}
\notetext{29}{既替余以蕙纕兮：既然因为我佩蕙草而罢免我。蕙（huì），香草；纕（xiāng），佩带。以佩香草为"罪"。}
\notetext{30}{又申之以揽茝：又加上我采摘白芷。申，加；茝（chǎi），白芷。摘取香草是美德，却被视为罪过——写尽正直之人在浊世中的冤屈。}
\notetext{31}{亦余心之所善兮：只要是我心中认为美好的。}
\notetext{32}{虽九死其犹未悔：即使死九次也不会后悔。九死，极言死之多次。"九死未悔"后来成为形容执着信念的成语。}
\notetext{33}{怨灵修之浩荡兮：怨恨灵修（楚王）糊涂放纵。浩荡，糊涂、放纵无度。}
\notetext{34}{终不察夫民心：始终不体察人心（我的忠心）。}
\notetext{35}{众女嫉余之蛾眉兮：那些女子嫉妒我的蛾眉。众女，比喻朝中的小人；蛾眉，像蚕蛾触角一样弯弯的细眉，比喻美好的品质。}
\notetext{36}{谣诼谓余以善淫：造谣诽谤说我善于淫邪。谣诼（zhuó），造谣诬蔑。}
\notetext{37}{固时俗之工巧兮：时俗本来就善于取巧。工巧，善于投机取巧。}
\notetext{38}{偭规矩而改错：背弃规矩，改易措施。偭（miǎn），违背；错，同"措"。}
\notetext{39}{背绳墨以追曲兮：背离绳墨（正直的标准）去追求邪曲。绳墨，木匠画直线的工具，比喻正道。}
\notetext{40}{竞周容以为度：争相周旋讨好，把这当作常规。周容，奉承献媚。}
\notetext{41}{忳郁邑余侘傺兮：我愁苦忧郁，失意不安。忳（tún），忧愁；郁邑，郁闷；侘傺（chà chì），失意的样子。}
\notetext{42}{吾独穷困乎此时也：我独自在这时候陷入穷困。}
\notetext{43}{宁溘死以流亡兮：宁愿突然死去，魂魄流散。溘（kè），忽然。}
\notetext{44}{余不忍为此态也：我也不忍做出那种（谄媚的）姿态。}
\notetext{45}{鸷鸟之不群兮：猛禽不与其他鸟合群。鸷（zhì）鸟，鹰隼之类的猛禽。}
\notetext{46}{自前世而固然：自古以来就是如此。}
\notetext{47}{何方圜之能周兮：方的和圆的怎么能合到一起？方圜（huán），方和圆；周，合。}
\notetext{48}{夫孰异道而相安：志向不同的人怎能相安无事？}
\notetext{49}{屈心而抑志兮：委屈内心，压抑志向。}
\notetext{50}{忍尤而攘诟：忍受罪责和耻辱。尤，罪责；攘（rǎng），容忍；诟（gòu），耻辱。}
\notetext{51}{伏清白以死直兮：坚守清白，为正直而死。伏，守。}
\notetext{52}{固前圣之所厚：这本来就是前代圣人所看重的。}
\notetext{53}{悔相道之不察兮：后悔当初选择道路时没有看清。相（xiàng），观察、选择。}
\notetext{54}{延伫乎吾将反：长久地站立，我打算回头了。延伫（zhù），久久站立；反，同"返"。}
\notetext{55}{回朕车以复路兮：调转我的车，回到原来的路上。}
\notetext{56}{及行迷之未远：趁着迷路还不太远。}
\notetext{57}{步余马于兰皋兮：让我的马在长着兰草的水边漫步。皋（gāo），水边的高地。}
\notetext{58}{驰椒丘且焉止息：奔驰到长着椒树的山丘上，暂且在那里歇息。椒丘，椒树丛生的小丘。}
\notetext{59}{进不入以离尤兮：进谏不被接纳，反而获罪。离，通"罹"，遭受。}
\notetext{60}{退将复修吾初服：退下来重新修整我当初的服饰（恢复本来的品格）。}
\notetext{61}{制芰荷以为衣兮：把菱叶和荷叶裁成上衣。芰（jì），菱角。}
\notetext{62}{集芙蓉以为裳：把荷花收集起来做成下裳。芙蓉，荷花。以植物为衣裳，象征高洁不染。}
\notetext{63}{不吾知其亦已兮：没人了解我也就算了吧。已，罢了。}
\notetext{64}{苟余情其信芳：只要我的内心确实芬芳。苟，只要；信，确实。}
\notetext{65}{高余冠之岌岌兮：把我的帽子做得高高的。岌岌（jí），高耸的样子。高冠是屈原的标志性装束，象征不与世俗同流。}
\notetext{66}{长余佩之陆离：把我的佩带做得长长的。陆离，长长的样子（一说光彩闪耀）。}
\notetext{67}{芳与泽其杂糅兮：芳香和润泽混杂在一起。}
\notetext{68}{唯昭质其犹未亏：只有光明纯洁的品质还没有亏损。}
\notetext{69}{忽反顾以游目兮：忽然回头四处张望。}
\notetext{70}{将往观乎四荒：将要去看看四方遥远的地方。四荒，四方边远之地。}
\notetext{71}{佩缤纷其繁饰兮：佩饰缤纷繁盛。}
\notetext{72}{芳菲菲其弥章：香气菲菲，更加浓郁彰明。弥，更加；章，同"彰"，明显。}
\notetext{73}{民生各有所乐兮：人生各有各的喜好。}
\notetext{74}{余独好修以为常：我独独爱好修身，习以为常。}
\notetext{75}{虽体解吾犹未变兮：即使身体被肢解，我也不改变。体解，肢解——古代最残酷的刑罚之一。}
\notetext{76}{岂余心之可惩：难道我的心会因为受到打击而改变吗？惩，受创而改变、屈服。全诗收束于一个反问——外力能伤害身体，但改变不了内心。}
\end{annotations}

\begin{appreciation}
《离骚》是中国文学史上最长的抒情诗，共三百七十三句、两千四百九十字。这里节选了三个段落：开篇自述身世，中间政论抒怀，以及最后的出处抉择。合在一起，便是一首完整的灵魂史诗。

诗的开篇，屈原就做出了一个惊人的宣告：他是帝高阳的后裔，生于寅年寅月寅日——这在当时被认为是极为吉祥的命格。他的父亲给了他"正则"和"灵均"两个名字，一个指向公正的法则，一个指向灵善的均衡。这个名字本身就是一种使命：他要做一个正直而美好的人。然后他写自己披香草、佩秋兰，朝夕采摘木兰宿莽——用一切美好的事物来滋养自己。但越是美好的人，在浊世中就越痛苦。

中间的政论部分，是全诗最激越的段落。"长太息以掩涕兮，哀民生之多艰"，屈原以泪眼回望自己走过的路。他早上在朝堂上进谏，晚上就被罢免——"謇朝谇而夕替"七个字，写尽了一个忠臣的宿命。但他没有求饶，没有收敛，反而说出了那句千古名言："亦余心之所善兮，虽九死其犹未悔。"只要是我心中认为对的，死九次也不后悔。这不是倔强，是一种对道义的绝对忠诚。

"怨灵修之浩荡兮，终不察夫民心"——楚王糊涂，始终不体察他的忠心。众女嫉余之蛾眉——朝廷中的小人们嫉妒他的才华，造谣诽谤。但屈原说：宁愿突然死去，也绝不做谄媚的姿态。"何方圜之能周兮，夫孰异道而相安？"方和圆不能合在一起，道不同的人不能相安。这是屈原最清醒的自觉：他不是没有想过妥协，而是知道——在这条路上，一旦妥协，就不再是自己了。

最后的段落写于行将远游之际。"悔相道之不察兮，延伫乎吾将反"，他后悔没有看清道路，打算掉头回去了。这是《离骚》中最柔软、最人性的时刻：承认自己的失败，承认自己可能选错了路。但他说的"回去"，不是放弃理想，而是回去修整"初服"——那个最初的、未受污染的自我。制芰荷以为衣，集芙蓉以为裳——他要用荷花重新做一套衣裳，从内而外地洁净自己。

"虽体解吾犹未变兮，岂余心之可惩"——这是全诗的最后一句，也是屈原给世界的最后答案。身体可以被肢解，但心不可被惩戒。你可以摧毁我，但你无法改变我。这句话的分量，要用屈原的人生来称量：他最后确实走向了汨罗江，用生命兑现了这句誓言。他死了，但他的心，两千三百年了，从来没有被惩戒过。
\end{appreciation}
